%----------------------------------------------------------------------------------------
% INVERSE, DETERMINANT EN RANG. Alle opgaven zijn eigen werk.
%
% Opbouw: Start, leerstof (theorie, werkwijze of uitleg, voorbeeld, korte oefening), en op een
% nieuwe pagina de Opdrachten (door elkaar, stijgende moeilijkheid), Koppel en Test jezelf.
% Labels van dit hoofdstuk beginnen met h3-.
%----------------------------------------------------------------------------------------

\mychapter{3}{Inverse, determinant en rang}

\lkr{Leerplan B+S: LPD 25-27 (matrices), met de wenken: de inverse via $[M \,|\, I] \to [I \,|\, M^{-1}]$, de determinant ingeleid met de inverse van een $2 \times 2$-matrix, ontwikkelen naar een rij of kolom en de regel van Sarrus, en het verband tussen rang, inverteerbaarheid, determinant en stelsels. De regel van Cramer (wenk bij LPD 37) staat bij de determinant. Grote berekeningen mogen met het rekentoestel.}

Een vergelijking $3x = 12$ los je op door beide leden met $\frac13$ te vermenigvuldigen, het omgekeerde van 3. Kan dat ook met een matrixvergelijking $AX = B$? Delen door een matrix bestaat niet, maar soms is er een matrix die de rol van $\frac13$ speelt: de inverse. In dit hoofdstuk leer je wanneer die bestaat en hoe je ze berekent. Een enkel getal, de determinant, zegt meteen of ze bestaat. En de rang telt hoeveel echt verschillende vergelijkingen er in een stelsel zitten.

\begin{oefeningen}[Start]
\oef{h3-start-1}

\oef{h3-verken}
\end{oefeningen}

\lkr{Laat oefening \ref{oef:h3-verken} maken vóór de uitleg. In a vinden leerlingen dat $AB = BA = I$. In b lossen ze een stelsel op door met $B$ te vermenigvuldigen, zonder rijherleiden. In c merken ze dat het voor sommige matrices niet lukt: de tweede rij is het dubbel van de eerste.}

%========================================================================================
\section{De inverse matrix}

\subsection{Definitie}

Bij getallen is $\frac13$ het omgekeerde van 3, omdat $3 \cdot \frac13 = 1$. Bij matrices speelt de eenheidsmatrix $I$ de rol van 1. Omdat het product van matrices niet commutatief is, eisen we dat het in beide volgordes lukt.

\begin{theorie}
Een vierkante matrix $A$ is \textbf{\textcolor{accent}{inverteerbaar}} als er een matrix $B$ bestaat met \quad $\formule{AB = BA = I}$

\vspace{1mm}
Die matrix heet de \textbf{\textcolor{accent}{inverse}} van $A$, en je noteert ze $\formule{A^{-1}}$.

\begin{stelling}{De inverse is uniek}
Een vierkante matrix heeft hoogstens één inverse.
\end{stelling}

\begin{bewijs}
Stel dat $B$ en $C$ allebei inversen zijn van $A$. Met de associativiteit van het product is
\bewijsregel{B = BI = B(AC) = (BA)C = IC = C .}
\end{bewijs}
\end{theorie}

Niet elke vierkante matrix is inverteerbaar, net zoals 0 geen omgekeerde heeft. In de Verken lukte het niet voor $\km{1 & 2 \\ 2 & 4}$. En in hoofdstuk 1 zag je matrices $A \neq O$ en $B \neq O$ met $AB = O$: zo'n $A$ heeft geen inverse, want anders was $B = A^{-1}AB = A^{-1}O = O$.

\begin{oefeningen}
\oef{h3-isinverse}
\end{oefeningen}

\subsection{Matrixvergelijkingen}

Is $A$ inverteerbaar, dan los je $AX = B$ op zoals $3x = 12$: vermenigvuldig beide leden met $A^{-1}$. Omdat de volgorde telt, vermenigvuldig je aan dezelfde kant als waar $A$ staat.

\begin{theorie}
Voor een inverteerbare matrix $A$: \qquad $\formule{AX = B \iff X = A^{-1}B}$ \qquad $\formule{XA = B \iff X = BA^{-1}}$
\end{theorie}

Uit $AX = B$ volgt $A^{-1}AX = A^{-1}B$, dus $IX = A^{-1}B$. Schrijf je $BA^{-1}$, dan klopt het niet: $A^{-1}$ moet links staan, net als $A$.

\voorbeeld{}{Los op: $\begin{cases} 2x + y = 4 \\ 5x + 3y = 11 \end{cases}$, wetende dat $\km{2 & 1 \\ 5 & 3}^{-1} = \km{3 & -1 \\ -5 & 2}$.\label{vb:h3-vgl}}{%
\uitleg{Het stelsel is $AX = B$. Omdat we $A^{-1}$ kennen, is de oplossing gewoon een product.}
\begin{stappen}
X = A^{-1}B = \begin{bmatrix} 3 & -1 \\ -5 & 2 \end{bmatrix} \begin{bmatrix} 4 \\ 11 \end{bmatrix} & Links vermenigvuldigen met $A^{-1}$. \\[3.5ex]
X = \begin{bmatrix} 12 - 11 \\ -20 + 22 \end{bmatrix} = \begin{bmatrix} 1 \\ 2 \end{bmatrix} & Dus $x = 1$ en $y = 2$. Controle: $2 + 2 = 4$ en $5 + 6 = 11$.
\end{stappen}}

Dat is vooral handig als je hetzelfde stelsel met veel verschillende rechterleden moet oplossen: je berekent $A^{-1}$ één keer.

\begin{oefeningen}
\oef{h3-vglkort}
\end{oefeningen}

\subsection{De inverse berekenen}

Hoe vind je $A^{-1}$? De kolommen van $A^{-1}$ zijn de oplossingen van drie stelsels: $AX = $ de eerste, tweede en derde kolom van $I$. Die stelsels hebben dezelfde coëfficiëntenmatrix, dus rijherleid je ze in één keer.

\begin{theorie}
\begin{stelling}{Inverse met Gauss-Jordan}
Een vierkante matrix $A$ is inverteerbaar als en slechts als $A$ rijherleidt tot $I$. Dan rijherleidt \quad $\formule{[A \,|\, I] \ \to \ [I \,|\, A^{-1}]}$
\end{stelling}
\end{theorie}

We bewijzen deze stelling niet. Krijg je links een nulrij, dan is $A$ niet inverteerbaar, en stop je.

\voorbeeld{}{Bereken de inverse van $A = \begin{bmatrix} 1 & 2 & 1 \\ 0 & 1 & 1 \\ 1 & 2 & 2 \end{bmatrix}$.\label{vb:h3-inverse}}{%
\uitleg{We zetten $I$ naast $A$ en rijherleiden tot links $I$ staat. Elke rijbewerking doen we op de hele rij, ook rechts van de streep.}
\begin{stappenbreed}
\um{ccc|ccc}{1 & 2 & 1 & 1 & 0 & 0 \\ 0 & 1 & 1 & 0 & 1 & 0 \\ 1 & 2 & 2 & 0 & 0 & 1} & De matrix $[A \,|\, I]$. \\[5ex]
\um{ccc|ccc}{1 & 0 & -1 & 1 & -2 & 0 \\ 0 & 1 & 1 & 0 & 1 & 0 \\ 0 & 0 & 1 & -1 & 0 & 1} \ \begin{array}{l} \rb{R_1 - 2R_2} \\ \\ \rb{R_3 - R_1} \end{array} & $R_3 - R_1$ met de oorspronkelijke $R_1$. \\[5ex]
\um{ccc|ccc}{1 & 0 & 0 & 0 & -2 & 1 \\ 0 & 1 & 0 & 1 & 1 & -1 \\ 0 & 0 & 1 & -1 & 0 & 1} \ \begin{array}{l} \rb{R_1 + R_3} \\ \rb{R_2 - R_3} \\ \\ \end{array} & Links staat $I$. \\[5ex]
A^{-1} = \begin{bmatrix} 0 & -2 & 1 \\ 1 & 1 & -1 \\ -1 & 0 & 1 \end{bmatrix} & Controle: de eerste rij van $A$ maal de eerste kolom van $A^{-1}$ is $0 + 2 - 1 = 1$, enzovoort.
\end{stappenbreed}}

\begin{werkwijze}
Schrijf $[A \,|\, I]$ en rijherleid zoals bij Gauss-Jordan. Staat links $I$, dan staat rechts $A^{-1}$. Krijg je links een nulrij, dan is $A$ niet inverteerbaar. Controleer met $AA^{-1} = I$, minstens voor één rij en kolom.
\end{werkwijze}

\begin{theorie}
\begin{stelling}{Rekenregels voor de inverse}
Voor inverteerbare $n \times n$-matrices $A$ en $B$:

\vspace{1mm}
\qquad $\formule{(A^{-1})^{-1} = A} \qquad \formule{(AB)^{-1} = B^{-1}A^{-1}}$
\end{stelling}

\begin{bewijs}
De eerste: $A^{-1}A = AA^{-1} = I$ zegt dat $A$ de inverse is van $A^{-1}$. De tweede: met de associativiteit is
\bewijsregel{(AB)(B^{-1}A^{-1}) = A(BB^{-1})A^{-1} = AIA^{-1} = AA^{-1} = I ,}
en op dezelfde manier $(B^{-1}A^{-1})(AB) = I$.
\end{bewijs}
\end{theorie}

Let op de volgorde in $(AB)^{-1} = B^{-1}A^{-1}$: zoals bij kousen en schoenen. Wie eerst kousen en dan schoenen aandoet, trekt eerst de schoenen uit.

\nujij{Bereken met de inverse uit voorbeeld \ref{vb:h3-inverse} de oplossing van $AX = \km{4 \\ 1 \\ 5}$.}{$X = A^{-1}\km{4 \\ 1 \\ 5} = \km{0 - 2 + 5 \\ 4 + 1 - 5 \\ -4 + 0 + 5} = \km{3 \\ 0 \\ 1}$.}

\begin{oefeningen}
\oef{h3-inverse}

\oef{h3-matrixvgl}
\end{oefeningen}

%========================================================================================
\section{De determinant}

\subsection{De determinant van een \texorpdfstring{$2 \times 2$}{2x2}-matrix}

Zoek je met $[A \,|\, I]$ de inverse van een algemene matrix $A = \km{a & b \\ c & d}$, dan komt telkens hetzelfde getal $ad - bc$ in de noemer. Dat getal beslist of de inverse bestaat.

\begin{theorie}
De \textbf{\textcolor{accent}{determinant}} van $A = \begin{bmatrix} a & b \\ c & d \end{bmatrix}$ is \quad $\formule{\det A = \begin{vmatrix} a & b \\ c & d \end{vmatrix} = ad - bc}$

\begin{stelling}{Inverse van een $2 \times 2$-matrix}
Is $\det A \neq 0$, dan is \quad $\formule{A^{-1} = \dfrac{1}{ad - bc} \begin{bmatrix} d & -b \\ -c & a \end{bmatrix}}$
\end{stelling}

\begin{bewijs}
Reken het product uit:
\bewijsregel{\begin{bmatrix} a & b \\ c & d \end{bmatrix} \begin{bmatrix} d & -b \\ -c & a \end{bmatrix} = \begin{bmatrix} ad - bc & -ab + ba \\ cd - dc & -cb + da \end{bmatrix} = (ad - bc) I ,}
dus na deling door $ad - bc$ krijg je $I$. In de andere volgorde gaat het net zo.
\end{bewijs}
\end{theorie}

De inverse van een $2 \times 2$-matrix schrijf je dus zo op: verwissel de getallen op de hoofddiagonaal, verander het teken van de andere twee, en deel door de determinant. Is $\det A = 0$, dan is $A$ niet inverteerbaar: dat bewijs je in opdracht \ref{oef:h3-o11}.

\voorbeeld{}{Bepaal de inverse van $A = \begin{bmatrix} 3 & 1 \\ 2 & -1 \end{bmatrix}$.\label{vb:h3-inverse2}}{%
\uitleg{We berekenen eerst de determinant. Is die niet 0, dan gebruiken we de formule.}
\begin{stappen}
\det A = 3 \cdot (-1) - 1 \cdot 2 = -5 & Niet 0: $A$ is inverteerbaar. \\[1ex]
A^{-1} = \frac{1}{-5}\begin{bmatrix} -1 & -1 \\ -2 & 3 \end{bmatrix} = \begin{bmatrix} \frac15 & \frac15 \\[.5ex] \frac25 & -\frac35 \end{bmatrix} & Diagonaal verwisselen, de andere twee van teken veranderen, delen door $-5$.
\end{stappen}}

\begin{oefeningen}
\oef{h3-inverse2k}
\end{oefeningen}

\subsection{De determinant van een \texorpdfstring{$3 \times 3$}{3x3}-matrix}

Ook bij grotere vierkante matrices bestaat een determinant, een getal dat zegt of de matrix inverteerbaar is. Je berekent ze door ze terug te brengen tot determinanten van $2 \times 2$-matrices.

\begin{theorie}
De \textbf{\textcolor{accent}{minor}} $M_{ij}$ van $A$ is de determinant van de matrix die overblijft als je rij $i$ en kolom $j$ schrapt. De determinant van een $3 \times 3$-matrix $A$, \textbf{\textcolor{accent}{ontwikkeld naar de eerste rij}}, is

\vspace{1mm}
\qquad $\formule{\det A = a_{11}M_{11} - a_{12}M_{12} + a_{13}M_{13}}$

\vspace{1mm}
Je mag ook naar een andere rij of kolom ontwikkelen, met de tekens volgens het schaakbordpatroon $\begin{smallmatrix} + & - & + \\ - & + & - \\ + & - & + \end{smallmatrix}$: de uitkomst is altijd dezelfde.
\end{theorie}

Dat elke rij en kolom hetzelfde geeft, nemen we aan zonder bewijs. Kies dus een rij of kolom met veel nullen: elke 0 is een term minder.

\voorbeeld{}{Bereken $\det B$ voor $B = \begin{bmatrix} 2 & -1 & 3 \\ 1 & 0 & 2 \\ 4 & 1 & -1 \end{bmatrix}$.\label{vb:h3-det}}{%
\uitleg{De tweede rij bevat een 0, dus ontwikkelen we naar die rij. De tekens in rij 2 zijn $-$, $+$, $-$.}
\begin{stappenbreed}
\det B = -1 \cdot \begin{vmatrix} -1 & 3 \\ 1 & -1 \end{vmatrix} + 0 - 2 \cdot \begin{vmatrix} 2 & -1 \\ 4 & 1 \end{vmatrix} & Schrap telkens de rij en de kolom van het element. \\[4ex]
= -1 \cdot (1 - 3) - 2 \cdot (2 + 4) = 2 - 12 = -10 & Twee determinanten van $2 \times 2$-matrices.
\end{stappenbreed}}

Bij een $3 \times 3$-matrix kan het ook met de \textbf{regel van Sarrus}. Schrijf de eerste twee kolommen nog eens rechts van de matrix. Tel de producten op de drie dalende diagonalen op, en trek die op de drie stijgende diagonalen af.

\naast{\begin{minipage}{9.4cm}
Voor de matrix $B$ uit het voorbeeld:
\[ \det B = (0 - 8 + 3) - (0 + 4 + 1) = -10 . \]
De dalende diagonalen geven $2 \cdot 0 \cdot (-1)$, $(-1) \cdot 2 \cdot 4$ en $3 \cdot 1 \cdot 1$, de stijgende $3 \cdot 0 \cdot 4$, $2 \cdot 2 \cdot 1$ en $(-1) \cdot 1 \cdot (-1)$. Sarrus werkt alleen bij $3 \times 3$-matrices.
\end{minipage}}{\begin{tikzpicture}[font=\small]
  \matrix[getallen,nodes={minimum width=6mm}] (m) {2 \& -1 \& 3 \& |[grijs]| 2 \& |[grijs]| -1 \\ 1 \& 0 \& 2 \& |[grijs]| 1 \& |[grijs]| 0 \\ 4 \& 1 \& -1 \& |[grijs]| 4 \& |[grijs]| 1 \\};
  \foreach \c in {1,2,3}{\pgfmathtruncatemacro{\e}{\c+2}\draw[accent,thick,opacity=.55] (m-1-\c.north west) -- (m-3-\e.south east);}
  \foreach \c in {3,4,5}{\pgfmathtruncatemacro{\b}{\c-2}\draw[fel!80!black,thick,dashed,opacity=.8] (m-1-\c.north east) -- (m-3-\b.south west);}
\end{tikzpicture}}

\begin{werkwijze}
Bij een $2 \times 2$-matrix: $ad - bc$. Bij een $3 \times 3$-matrix: ontwikkel naar een rij of kolom met veel nullen, met de tekens van het schaakbord, of gebruik de regel van Sarrus. Controleer met het rekentoestel als het mag.
\end{werkwijze}

\begin{oefeningen}
\oef{h3-det}
\end{oefeningen}

\subsection{Eigenschappen}

\begin{theorie}
\begin{stelling}{Eigenschappen van de determinant}
Voor vierkante $n \times n$-matrices $A$ en $B$:

\vspace{1mm}
{\centering\arrayrulecolor{accent}\renewcommand{\arraystretch}{1.5}%
\begin{tabular}{|>{\centering\arraybackslash}m{6.4cm}|>{\centering\arraybackslash}m{5.6cm}|}
\hline
product & $\formule{\det(AB) = \det A \cdot \det B}$ \\ \hline
getransponeerde & $\formule{\det A\tr = \det A}$ \\ \hline
twee rijen verwisselen & de determinant verandert van teken \\ \hline
een rij vermenigvuldigen met $r$ & de determinant wordt $r$ keer zo groot \\ \hline
een veelvoud van een rij bij een andere optellen & de determinant verandert niet \\ \hline
\end{tabular}\par}
\end{stelling}

\begin{stelling}{Inverteerbaar}
Een vierkante matrix $A$ is inverteerbaar als en slechts als $\det A \neq 0$. Dan is \quad $\formule{\det(A^{-1}) = \dfrac{1}{\det A}}$
\end{stelling}
\end{theorie}

Die eigenschappen bewijzen we niet in het algemeen; voor $2 \times 2$-matrices kun je ze nagaan. Ze verklaren wel de stelling: rijherleiden verandert de determinant hoogstens met een factor verschillend van 0. Een matrix die tot $I$ rijherleidt, heeft dus een determinant verschillend van 0, en een matrix met een nulrij in de rijgereduceerde vorm heeft determinant 0. Zo heeft een matrix met twee gelijke rijen determinant 0: trek de ene af van de andere en je krijgt een nulrij.

\begin{oefeningen}
\oef{h3-eigensch}
\end{oefeningen}

\subsection{De regel van Cramer}

Bij een stelsel met evenveel vergelijkingen als onbekenden en een determinant verschillend van 0, kun je de oplossing ook met determinanten opschrijven.

\begin{theorie}
\begin{stelling}{Regel van Cramer}
Is $\det A \neq 0$, dan heeft $AX = B$ precies één oplossing, met \quad $\formule{x_j = \dfrac{\det A_j}{\det A}}$

\vspace{1mm}
Hierin is $A_j$ de matrix $A$ waarin kolom $j$ vervangen is door $B$.
\end{stelling}
\end{theorie}

\voorbeeld{}{Los op met de regel van Cramer: $\begin{cases} 3x - 2y = 4 \\ 5x + y = 11 \end{cases}$\label{vb:h3-cramer}}{%
\uitleg{We berekenen de determinant van $A$, en daarna die van $A$ met eerst de eerste, dan de tweede kolom vervangen door de rechterleden.}
\begin{stappen}
\det A = \begin{vmatrix} 3 & -2 \\ 5 & 1 \end{vmatrix} = 3 + 10 = 13 & Niet 0: precies één oplossing. \\[3.5ex]
x = \frac{1}{13}\begin{vmatrix} 4 & -2 \\ 11 & 1 \end{vmatrix} = \frac{4 + 22}{13} = 2 & Kolom 1 vervangen door $\km{4 \\ 11}$. \\[3.5ex]
y = \frac{1}{13}\begin{vmatrix} 3 & 4 \\ 5 & 11 \end{vmatrix} = \frac{33 - 20}{13} = 1 & Kolom 2 vervangen.
\end{stappen}}

Voor grote stelsels is Cramer veel meer rekenwerk dan Gauss-Jordan. De regel is vooral handig bij twee onbekenden, of als je maar één onbekende nodig hebt.

\begin{oefeningen}
\oef{h3-cramer}
\end{oefeningen}

%========================================================================================
\section{De rang}

In een stelsel kan een vergelijking een gevolg zijn van de andere, zoals in voorbeeld \ref{vb:h2-onbepaald} van hoofdstuk 2: na rijherleiden werd ze een nulrij. De rang telt hoeveel vergelijkingen er overblijven.

\begin{theorie}
De \textbf{\textcolor{accent}{rang}} van een matrix $A$ is het aantal spillen in de rijgereduceerde vorm van $A$. Je noteert $\formule{\operatorname{rang} A}$.
\end{theorie}

Het aantal spillen is ook het aantal rijen dat niet nul is in de rijgereduceerde vorm. Omdat elke spil in een eigen rij en een eigen kolom staat, is de rang van een $m \times n$-matrix hoogstens het kleinste van $m$ en $n$.

\voorbeeld{}{Bepaal de rang van $R = \begin{bmatrix} 1 & 2 & -1 & 3 \\ 2 & 4 & 0 & 2 \\ 3 & 6 & -1 & 5 \end{bmatrix}$.\label{vb:h3-rang}}{%
\uitleg{We rijherleiden en tellen de spillen.}
\begin{stappenbreed}
\begin{bmatrix} 1 & 2 & -1 & 3 \\ 0 & 0 & 2 & -4 \\ 0 & 0 & 2 & -4 \end{bmatrix} \ \begin{array}{l} \\ \rb{R_2 - 2R_1} \\ \rb{R_3 - 3R_1} \end{array} & Kolom 2 krijgt geen spil: onder de 2 staan alleen nullen. \\[5ex]
\begin{bmatrix} 1 & 2 & 0 & 1 \\ 0 & 0 & 1 & -2 \\ 0 & 0 & 0 & 0 \end{bmatrix} \ \begin{array}{l} \rb{R_1 + \frac12 R_2} \\ \rb{\frac12 R_2} \\ \rb{R_3 - R_2} \end{array} & Twee spillen: $\operatorname{rang} R = 2$.
\end{stappenbreed}
De derde rij van $R$ is de som van de eerste twee: ze brengt niets nieuws.}

Met de rang vat je alles samen wat je over stelsels en inverteerbaarheid weet.

\begin{theorie}
\begin{stelling}{Rang, stelsels en inverteerbaarheid}
Een stelsel $AX = B$ met $n$ onbekenden is

\vspace{1mm}
\quad strijdig als en slechts als $\formule{\operatorname{rang} A < \operatorname{rang}\,[A \,|\, B]}$;

\quad anders heeft het $\formule{n - \operatorname{rang} A}$ parameters (bepaald als dat 0 is).

\vspace{1mm}
Voor een vierkante $n \times n$-matrix $A$ zijn gelijkwaardig: \ $\operatorname{rang} A = n$ \ $\iff$ \ $A$ is inverteerbaar \ $\iff$ \ $\det A \neq 0$ \ $\iff$ \ $AX = B$ is bepaald voor elke $B$.
\end{stelling}
\end{theorie}

De eerste regel is de strijdige rij $[0 \ \cdots \ 0 \,|\, c]$ uit hoofdstuk 2: die geeft $[A \,|\, B]$ een spil meer dan $A$. De tweede regel zegt dat elke onbekende zonder spil een parameter is.

\nujij{Een stelsel met 4 onbekenden heeft een coëfficiëntenmatrix met rang 3, en de uitgebreide matrix heeft ook rang 3. Hoeveel oplossingen zijn er?}{Niet strijdig, en $4 - 3 = 1$ parameter: oneindig veel oplossingen.}

\begin{oefeningen}
\oef{h3-rang}

\oef{h3-rangst}
\end{oefeningen}

%========================================================================================
\opdrachtenkop

\begin{oefeningen}
\needspace{14\baselineskip}
\oef{h3-o14}

\oef{h3-o17}

\oef{h3-o4}

\needspace{14\baselineskip}
\oef{h3-o15}

\oef{h3-o19}

\oef{h3-o2}

\oef{h3-o16}

\needspace{14\baselineskip}
\oef{h3-o18}

\oef{h3-o1}

\oef{h3-o3}

\needspace{14\baselineskip}
\oef{h3-o21}

\oef{h3-o5}

\needspace{14\baselineskip}
\oef{h3-o24}

\oef{h3-o7}

\oef{h3-o8}

\needspace{14\baselineskip}
\oef{h3-o20}

\oef{h3-o22}

\oef{h3-o6}

\oef{h3-o9}

\oef{h3-o23}

\oef{h3-o10}

\oef{h3-o25}

\oef{h3-o12}

\oef{h3-o13}

\oef{h3-o26}

\oef{h3-o11}
\end{oefeningen}

%----------------------------------------------------------------------------------------
% KOPPEL: verbanden met andere leerstof
%----------------------------------------------------------------------------------------
\needspace{20\baselineskip}
\opdrachtsectie{Koppel}

\lkr{Coördinaten en de oppervlakte van een driehoek kennen de leerlingen uit de 2de graad.}

\begin{oefeningen}
\oef{h3-koppel-1}

\end{oefeningen}

%----------------------------------------------------------------------------------------
\opdrachtsectie{Test jezelf}

\begin{oefeningen}
\oef{h3-t1}

\oef{h3-t2}

\oef{h3-t3}

\oef{h3-t4}

\oef{h3-t5}

\oef{h3-t6}

\oef{h3-t7}

\oef{h3-t8}
\end{oefeningen}

\herhaalschema{%
  \fv{oef:h3-t1}{vb:h3-inverse2}{oef:h3-inverse}%
  \fv{oef:h3-t2}{vb:h3-inverse}{oef:h3-inverse}%
  \fv{oef:h3-t3}{vb:h3-det}{oef:h3-det}%
  \fv{oef:h3-t4}{vb:h3-cramer}{oef:h3-cramer}%
  \fv{oef:h3-t5}{vb:h3-rang}{oef:h3-rang}%
  \fv{oef:h3-t6}{vb:h3-inverse2}{oef:h3-det}%
  \fv{oef:h3-t7}{vb:h3-vgl}{oef:h3-matrixvgl}%
  \fv{oef:h3-t8}{vb:h3-rang}{oef:h3-rangst}}

\ifzelfstudie
\section*{Maak je eigen samenvatting}
Sluit dit boek en schrijf in je schrift: wat een inverse matrix is en waarom ze uniek is, hoe je een matrixvergelijking oplost, hoe je de inverse berekent met $[A \,|\, I]$, de determinant en de inverse van een $2 \times 2$-matrix, hoe je een determinant van een $3 \times 3$-matrix berekent, de eigenschappen van de determinant, de regel van Cramer, wat de rang is, en het verband tussen rang, determinant, inverse en stelsels. Vergelijk daarna met de kaders.
\fi
